\section{The two preparatory flights}
\label{sec:preparation}

The first release has normal velocity of order $\eps$, and its flight
ends on the scale $\sqrt\eps$. The ensuing contact amplifies the
velocity to order $\eps^{3/4}$, so the next flight is longer, of
order $\eps^{3/8}$. Both flights inherit a nonzero normal memory.
We retain it in the initial opening layer and identify each actual
first touchdown before normalizing its complete incoming state.

Throughout $r=k\sqrt\eps=-\eta$ is the same physical belt offset
and $u_r=u_0-r$ is the physical slip. The belt coefficient $k$,
the remaining parameters and the initial preparation stay fixed.

The scales established below distinguish the first signed contact
from the recurring returns:
\[
\begin{array}{c|c|c}
 \text{touchdown}&\text{incoming scale}&\text{offset ratio}\\ \hline
 \TD_1&\sqrt\eps&r/\sqrt\eps=k\\
 \TD_2&\delta_2\asymp\eps^{3/8}&r/\delta_2^{4/3}\asymp1\\
 \TD_3&\delta_3\asymp\eps^{9/32}&r/\delta_3^{4/3}=O(\eps^{1/8})
\end{array}
\]
The second touchdown has a bounded complete cone. One further actual
return makes the offset ratio small and places the trajectory in the
domain used for iteration.

In a flight,
\begin{equation}
 \dot y=v,\quad\dot v=b(X),\quad\dot z=-z/\eps,
 \quad\dot u_r=a_T(X),\qquad g=y-z.
 \label{prep:free-field}
\end{equation}

\begin{lemma}[The gap of a local free extension]
\label{lem:detuned-free-gap-expansion}
Suppose $|X_L-G|\le C\rho$, $\rho\to0$, and $y_L=z_L$.
On a common fixed local interval, the smooth free extension satisfies
\begin{align}
 p(t)&=p_L-at+O(\rho t+t^2),\qquad
 b(t)=b_L-ct+O(\rho t+t^2),\nonumber\\
 u_r(t)&=u_{r,L}+d_\lambda t+O(\rho t+t^2),\nonumber\\
 v(t)&=v_L+b_Lt-ct^2/2+O(\rho t^2+t^3),
 \label{prep:free-expansions}\\
 g(t)&=z_L(1-e^{-t/\eps})+v_Lt+b_Lt^2/2-ct^3/6+\mathcal R(t),
 \nonumber\\
 |\mathcal R(t)|&\le C(\rho t^3+t^4),\qquad
 |\mathcal R'(t)|\le C(\rho t^2+t^3).
 \label{prep:gap-expansion}
\end{align}
If $\min\{v_L,v_L+z_L/\eps\}\ge v_*>0$, then
\begin{equation}
 \dot g(t)\ge v_*-C\rho t-Ct^2.
 \label{prep:opening}
\end{equation}
The extension is the actual unilateral solution on every initial
interval on which $g>0$.
At any positive root $T$ of this free gap within the local interval,
put $D=-p(T)/a$ and $H_b=\dot b+c$. Then
\begin{equation}
 -v(T)-\frac c3D^2=\frac{y_L-z(T)}T-\frac12b_LT+
 \frac c3(T^2-D^2)-\frac1{2T}\int_0^T(T^2-t^2)H_b(t)\,dt.
 \label{prep:touchdown-cancellation}
\end{equation}
\end{lemma}
\begin{proof}
Stop at a common chart boundary. The bounded free field gives
$|X(t)-G|\le C(\rho+t)$, which improves the stop on a smaller
fixed time interval. Its critical rates give
$\dot p=-a+O(\rho+t)$, $\dot b=-c+O(\rho+t)$ and
$a_T=d_\lambda+O(\rho+t)$. Successive integrations give all the
mechanical estimates and the gap remainder. The memory is exactly
$z_Le^{-t/\eps}$. Moreover
$v_L+(z_L/\eps)e^{-t/\eps}\ge\min(v_L,v_L+z_L/\eps)$,
which proves \eqref{prep:opening}. While $g>0$ both contact forces
vanish; uniqueness identifies the extension with the actual free leg.
For the last identity integrate $\dot b=-c+H_b$ twice:
\[
 \begin{aligned}
 v(T)&=v_L+b_LT-cT^2/2+\int_0^T(T-t)H_b(t)\,dt,\\
 z(T)&=y_L+v_LT+b_LT^2/2-cT^3/6
                  +\tfrac12\int_0^T(T-t)^2H_b(t)\,dt.
 \end{aligned}
\]
Eliminate $v_L$ with the second equality, which is the actual root
condition $y(T)=z(T)$. No asymptotic estimate is used in this identity.
\end{proof}

\subsection{The first touchdown and signed compression}

Use the critical-release data of Theorem~\ref{thm:detuned-release}
and put $s=\sqrt\eps$. Define
\begin{equation}
 H=ak/\gamma,\qquad q=q_{\rm in},\qquad
 T_f=\frac{\sqrt{9\ell^2H^2+24cq}-3\ell H}{2c},\qquad
 \pi=H-aT_f.
 \label{prep:first-flight-data}
\end{equation}
Thus $cT_f^2+3\ell HT_f=6q$. These quantities have common finite
bounds and $T_f,q,H$ have common positive lower bounds on the
fixed compact label classes.

\begin{proposition}[First flight and its amplified release]
\label{prep:first-flight}
The first free leg after the captured passage has a transverse first
touchdown, at elapsed time $t_1$, with
\begin{gather}
 t_1/s\to T_f,\qquad p_1/s\to\pi,\qquad
 b_1/s\to-\ell H-cT_f,\qquad u_{r,1}/s\to S:=d_\lambda T_f,\\
 -v_1/\eps\to Q:=T_f(\ell H/2+cT_f/3)>0,\qquad
 y_1=z_1=o(\eps^2),\qquad |X_1-G|=O(s).
 \label{prep:first-touchdown}
\end{gather}
The subsequent contact has its first release with
\begin{gather}
 |X_L-G|=O(s),\qquad p_L,b_L=O(s),\qquad u_{r,L}=O(\eps),\\
 v_L=\eps^{3/4}(R+o(1)),\qquad
 y_L=z_L=\eps^{7/4}(Z+o(1))<0,
 \label{prep:amplified-release}\\
 R=A_s\nu_\#,\quad Z=-A_s\gamma_\#,\qquad
 A_s^2=QS/h-\pi S^2/(2h^2)>0.
 \label{prep:amplitude}
\end{gather}
In particular $R$ and $R+Z$ have common positive lower bounds.
The first touchdown parameter $\pi$ may be positive, zero or negative.
\end{proposition}
\begin{proof}
The release theorem gives $p_L/s\to H$, $b_L/s\to-\ell H$,
$v_L/\eps\to q$ and a common positive opening margin
$(v_L+z_L/\eps)/\eps$. Apply the free-gap lemma with $\rho=s$.
For $t=s\tau$, uniformly on bounded intervals,
\[
 s^{-3}g(s\tau)\longrightarrow
 P_f(\tau):=q\tau-\ell H\tau^2/2-c\tau^3/6.
\]
The normalized memory value is $O(s)$; away from $\tau=0$ its
normalized derivative tends to zero exponentially. The opening
estimate gives $\dot g>0$ on $0\le t\le s_0s$ for a common small
$s_0>0$. The polynomial is positive until its unique positive zero
$T_f$, where $P_f'(T_f)=-T_f(\ell H/2+cT_f/3)<0$.
Choose a fixed small bracket with a common negative derivative
margin. Uniform value convergence preserves positivity between the
opening interval and the bracket and the opposite endpoint signs;
uniform derivative convergence in the bracket gives exactly one zero.
It is the actual first transverse touchdown. The mechanical expansions
give \eqref{prep:first-touchdown}; exact memory decay gives, more
precisely, $|z_1|\le C\eps^2e^{-c/\sqrt\eps}$ and makes
$\dot g(t_1)$ have the same negative leading term as $v_1$.

For the ensuing contact set $\delta=s$. Its exact normalized incoming
data are $\pi_\eps=p_1/\delta$, $Q_\eps=-v_1/\delta^2$,
$S_\eps=u_{r,1}/\delta$ and
$\zeta_\eps=z_1/(\eps\delta^2)$.
They converge to $(\pi,Q,S,0)$, and
\[
 2hQ-\pi S
 =T_f\{(h\ell-d_\lambda)H+(2hc/3+ad_\lambda)T_f\}>0,
\]
with a common strict margin since $h\ell-d_\lambda=\gamma>0$.
Thus they lie in a fixed compact subset of the signed contact domain
of Theorem~\ref{con:contact-release}; the incoming gap-closing margin
also holds. Here $\eps/\delta^2=1$ and $r=k\delta$.
That theorem, applied to this actual state, proves the first release,
the complete normal memory and \eqref{prep:amplified-release}.
Its normal reference has $\nu_\#>\gamma_\#>0$, giving the
opening margin $R+Z>0$. The strict amplitude inequality and compactness
bound $A_s$ away from zero. No new state has been prepared.
\end{proof}

At this first touchdown $r=k\delta$, so the signed contact theorem
supplies the amplified release. The next flight brings the offset
into the range $r\le A\delta^{4/3}$ required for recurring returns.

\subsection{The second touchdown in complete coordinates}

The normal displacement at the amplified release is too large to use
the first-flight normalization. Its exact decay still ensures an
initial free interval; on the longer flight the cubic acceleration
term dominates the initial $O(\sqrt\eps)$ acceleration. At touchdown,
the smaller fibre scale requires cancellation of the actual release
velocity before division.

\begin{lemma}[Complete preparation at the second touchdown]
\label{lem:td2-complete-preparation}
Suppose an actual release has $|X_L-G|=O(\sqrt\eps)$,
$u_{r,L}=O(\eps)$ and satisfies \eqref{prep:amplified-release},
with bounded $R,Z$ and common positive lower bounds for $R,R+Z$.
Its next touchdown is the transverse first positive zero of the free
gap. If $T$ is its elapsed time, $D=-p(T)/a$ and
$C_R=\sqrt{6R/c}$, then
\begin{gather}
 T=\eps^{3/8}C_R(1+o(1)),\qquad D\sim T,\qquad
 D-T=O(\sqrt\eps),\qquad v(T)/\eps^{3/4}\to-2R,
 \label{prep:second-scale}\\
 \eps/D^2=O(\eps^{1/4})\to0.
\end{gather}
Put $r_0=c/3$ and define the complete actual incoming coordinates
\begin{gather}
 U_2^r=\frac{u_r(T)}{d_\lambda D},\quad B_2=\frac{-b(T)}{cD},\quad
 W_2=\frac{-v(T)}{r_0D^2},\quad
 \zeta_2=\frac{z(T)}{\eps D^2},\\
 \eta_{U,2}^r=\frac{U_2^r-1}{D^{1/3}},\quad
 \eta_{B,2}=\frac{B_2-1}{D^{1/3}},\quad
 \eta_{W,2}=\frac{W_2-1}{D^{1/3}},\qquad
 \xi_2^r=(\zeta_2,\eta_{U,2}^r,\eta_{W,2}).
 \label{prep:second-coordinates}
\end{gather}
There exist fixed constants $M_2<\infty$ and $\gamma_2>0$ such that
\begin{gather}
 |\xi_2^r|+|\eta_{B,2}|\le M_2,\qquad
 r_0W_2-\zeta_2\ge\gamma_2,\qquad
 (B_2,U_2^r,W_2)\to(1,1,1),\\
 -Ce^{-c_1\eps^{-5/8}}\le\zeta_2<0,\qquad
 \frac r{D^{4/3}}=\frac{k}{C_R^{4/3}}(1+o(1)).
 \label{prep:complete-second-cone}
\end{gather}
The offset ratio has positive finite common upper and lower bounds.
\end{lemma}
\begin{proof}
Put $P=\sqrt\eps$, $L=\eps^{3/8}$, so the release velocity scale
is $L^2$ and its memory scale is $\eps L^2$. For $t=L\tau$,
the free-gap lemma gives
\[
 L^{-3}g(L\tau)\longrightarrow R\tau-c\tau^3/6.
\]
The initial quadratic term is $O(P/L)$, memory value is
$O(\eps/L)$ and the remainder is $O(P+L)$, all tending to zero.
Away from zero the normalized memory derivative decays exponentially.
On $0\le t\le s_0L$, \eqref{prep:opening} gives
$\dot g\ge L^2\{m-Cs_0P/L-Cs_0^2\}>0$ after choosing $s_0$
small and then $\eps$ small. The cubic has its unique positive zero
$C_R$ with derivative $-2R$. Its positive interval, opposite root
bracket signs and negative derivative margin transfer to the extension.
Thus the bracket contains precisely its first zero and is an actual
touchdown, with $T/L\to C_R$ and $v(T)/L^2\to-2R$.

All finer estimates are on free prefixes. Stop at a chart boundary
or $HL$, where $H>\sup C_R+1$ is fixed. The bounded free field
gives $|X(t)-G|\le C(P+t)$, improving the chart stop; hence
\begin{equation}
 A_p:=\dot p+a=O(P+t),\quad H_b:=\dot b+c=O(P+t),\quad
 \dot u_r-d_\lambda=O(P+t),\qquad 0\le t\le HL.
 \label{prep:prefix-residuals}
\end{equation}
In particular $D-T=-p_L/a-a^{-1}\int_0^TA_p\,dt=O(P)$, since
$T\asymp L$ and $L^2=o(P)$. This proves $D\sim T>0$.
Retain the critical combinations
\begin{equation}
 J_U=u_r+(d_\lambda/a)p,\qquad J_B=-b+(c/a)p,\qquad
 J_W=-v-(c/3)(p/a)^2.
 \label{prep:critical-combinations}
\end{equation}
The free-field derivatives of $J_U$ and $J_B$ vanish at $G$, so integration with
\eqref{prep:prefix-residuals} gives
$J_U(T),J_B(T)=O(P+PT+T^2)=O(P)$.

For $J_W=-v(T)-(c/3)D^2$, use
\eqref{prep:touchdown-cancellation}, which eliminates the release
velocity through the actual touchdown identity $y(T)=z(T)$.
The four terms have sizes $O(\eps^{11/8})$,
$O(\eps^{7/8})$, $O(\eps^{7/8})$ and $O(\eps^{9/8})$.
For the last use the bound on every prefix:
$C(PT^2+T^3)=O(\eps^{9/8})$. Therefore $J_W(T)=O(PL)$.
No rate for the $o(1)$ in the release velocity is required.
Only now divide, using the exact identities
\[
 \eta_{U,2}^r=\frac{J_U(T)}{d_\lambda D^{4/3}},\quad
 \eta_{B,2}=\frac{J_B(T)}{cD^{4/3}},\quad
 \eta_{W,2}=\frac{J_W(T)}{r_0D^{7/3}}.
\]
Their denominators are respectively of sizes $P,P,PL$.
The three fibres are bounded and the three unscaled ratios tend to
one. Finally $z(T)=z_Le^{-T/\eps}<0$ and
$|z_L|/(\eps D^2)\le C$; since
$T/\eps\ge c_1\eps^{-5/8}$, this proves the memory bound and
the strict closing margin. The last offset identity follows from
$D\sim C_R\eps^{3/8}$ and $r=k\sqrt\eps$.
\end{proof}

\begin{corollary}[The physical second touchdown]
\label{cor:original-to-td2}
The original physical trajectory has this complete second-touchdown
state, with $\delta_2=D\asymp\eps^{3/8}$ and $p_2<0$.
All bounds are uniform on its fixed compact label class, including
labels for which the leading first-touchdown parameter $\pi$ vanishes.
\end{corollary}
\begin{proof}
Proposition~\ref{prep:first-flight} supplies every release hypothesis
of Lemma~\ref{lem:td2-complete-preparation}, including $u_{r,L}=O(\eps)$,
negative memory, and the common $R,R+Z$ margins. The lemma applies
to that same trajectory.
\end{proof}

\begin{corollary}[Entry to the recurring returns]
\label{cor:td3-admission}
Write $\eta_{U,*,\lambda}$ for the slip component of
$\xi_{*,\lambda}$ and put
$\eta_{B,*,\lambda}=(g_b/c)\eta_{U,*,\lambda}$, with
$g_b=-d_\lambda b_u(G)$ as in Lemma~\ref{ret:inverse}.
The next actual contact--flight return reaches $\TD_3$ with
\begin{gather}
 \delta_3\asymp\eps^{9/32},\qquad
 \xi_3^r\in\operatorname{int}\Kcal,\qquad
 r/\delta_3^{4/3}=O(\eps^{1/8}),\\
 \left|\log\frac{\delta_3}{K_\lambda\delta_2^{3/4}}\right|
 +|\xi_3^r-\xi_{*,\lambda}|+
 |\eta_{B,3}-\eta_{B,*,\lambda}|
 \le C(\eps^{3/32}+\eps^{7/16}+\eps^{1/8}).
 \label{prep:third-entry-error}
\end{gather}
It retains the complete negative normal memory and every strict
physical margin of Theorem~\ref{thm:detuned-output-admission}, and is an admissible
input to a further return. In particular
$\log(1/\delta_3)=(9/32)\log(1/\eps)+O(1)$ uniformly.
\end{corollary}
\begin{proof}
Choose a fixed finite cone containing the full second-touchdown
bounds, and fixed $A$ covering $r/\delta_2^{4/3}$.
Since $\eps/\delta_2^2=O(\eps^{1/4})\to0$ and $\delta_2\to0$,
Theorem~\ref{thm:detuned-output-admission} applies to this actual input for all
sufficiently small $\eps$. Its three error terms are
$\delta_2^{1/4}=O(\eps^{3/32})$,
$\eps\delta_2^{-3/2}=O(\eps^{7/16})$, and
$r/\delta_2=O(\eps^{1/8})$.
It gives $\delta_3\asymp\delta_2^{3/4}$ before output admission.
Lemma~\ref{ret:inverse} recovers the dependent acceleration from the
actual output. Its additional terms are
$O(r/\delta_3^{4/3})=O(\eps^{1/8})$ and
$O(\delta_3^{2/3})=O(\eps^{3/16})$, within the displayed bound.
The strict output cone and memory are supplied by that theorem.
Since $\delta_3\to0$, the output eventually also satisfies its input
cutoff, permitting the next actual return. The uniform positive upper
and lower scale bounds prove the logarithmic statement, with the
original labels and belt parameter on the same physical trajectory.
\end{proof}
